\subsection{Commutative diagrams}

Consider, for example, the following diagram:

\[
\begin{array}{c} \mathrm{A} \xrightarrow {f} \mathrm{B} \xrightarrow {g} \mathrm{C} \xrightarrow {h} \mathrm{D} \\ a \Big \downarrow \qquad b \Big \downarrow \qquad c \Big \downarrow \qquad d \Big \downarrow \\ \mathrm{A} ^ {\prime} \xrightarrow [ f ^ {\prime} ]{} \mathrm{B} ^ {\prime} \xrightarrow [ g ^ {\prime} ]{} \mathrm{C} ^ {\prime} \xrightarrow [ h ^ {\prime} ]{} \mathrm{D} ^ {\prime} \end{array}\tag{2}
\]

To every path composed of a certain number of segments of the diagram traversed in the direction shown by the arrows corresponds a mapping of the set represented by the beginning of the first segment to the set represented by the end of the last segment, namely the composition of the mappings represented by the various segments traversed. For each vertex of the diagram, for example B, by way of convention there is a path reduced to B, to which corresponds the identity mapping \(l_{B}\) .

In (2) there are for example three paths beginning at A and ending at C'; the corresponding mappings are \(c \circ g \circ f\) , \(g' \circ b \circ f\) and \(g' \circ f' \circ a\) . A diagram is said to be commutative if, for every pair of paths in the diagram with the same beginning and end, the two corresponding mappings are equal; in particular if the beginning and end of a path coincide the corresponding mapping must be the identity.

For the diagram (2) to be commutative it is necessary and sufficient that the relations

\[
f ^ {\prime} \circ a = b \circ f, \quad g ^ {\prime} \circ b = c \circ g, \quad h ^ {\prime} \circ f = d \circ h;\tag{3}
\]

hold; in other words, it is necessary and sufficient that the three square diagrams contained in (2) be commutative. For the relations (3) imply \(c \circ g \circ f = g' \circ b \circ f\) since \(c \circ g = g' \circ b\) , and \(g' \circ b \circ f = g' \circ f' \circ a\) since \(b \circ f = f' \circ a\) ; thus the three paths beginning at A and ending at C' give the same mapping. It can be similarly verified that the four paths beginning at A and ending at D' (resp. the three paths beginning at B and ending at D') give the same mapping. The relations (3) signify that the two paths beginning at A (resp. B, C) and ending at B' (resp. C', D') give the same mapping. None of the other pairs of vertices of (2) can be joined by more than one path and the diagram (2) is therefore commutative.

In what follows we shall leave the reader to verify analogous results for other types of diagrams.

